\section{Full multiplicity from a restricted coefficient box}

\subsection*{Finding F-015: shape-sensitive padding preserves every copy}
Let \(\lambda\vdash nd\) have \(2\leq r\leq\min\{d,n^2\}\)
rows, and put \(b=\lambda_2\). For \(j\geq0\), define
\[
 C_j(\lambda)=\#\{\alpha\in\mathbb Z_{\geq0}^{r-1}:
       \alpha_i\leq\lambda_i\ (2\leq i\leq r),\quad
       |\alpha|\leq j\}.
\]
Set
\[
 \begin{split}
 u&=\left\lfloor\frac{b-1}{2}\right\rfloor,\qquad
 v=\left\lceil\frac{b-1}{2}\right\rceil,\\
 D_{\mathrm{box}}(\lambda)&=C_u(\lambda)+C_v(\lambda)-1.
 \end{split}
\]
Also let \(T_b(\lambda)\) denote the bounded variable-splitting
count of M-009, with its degree parameter replaced by \(b\):
\[
 \begin{split}
 N_A(b,\lambda)&=
 \#\{\alpha\in\mathbb Z_{\geq0}^{A}:\alpha_i\leq\lambda_i,
                                            |\alpha|\leq b\},\\
 T_b(\lambda)&=
 \min_{A\sqcup B=\{2,\ldots,r\}}
                  \bigl(N_A(b,\lambda)+N_B(b,\lambda)-1\bigr).
 \end{split}
\]
Then
\begin{equation}\label{eq:F015}
 n\geq\min\{D_{\mathrm{box}}(\lambda),T_b(\lambda)\}
 \quad\Longrightarrow\quad
                       K_n(\lambda)=a_\lambda(d[n]).
\end{equation}
The assertion holds at every outer degree \(d\). It excludes
actual multiplicity obstructions for every orbit closure in the
same ambient form space, including the internally padded permanent.
It is a full-multiplicity conclusion, stronger than mere occurrence.

The count \(D_{\mathrm{box}}\) is at most \(B(b,r)\) from F-014,
because imposing the coordinate bounds only deletes monomials.
When \(b\leq m\), the other count satisfies
\(T_b(\lambda)\leq T_m(\lambda)\). Thus the criterion can
strengthen both the uncapped second-row bound and M-009.

\begin{proof}
The coordinate bound on \(\alpha_2\) is \(b\). Hence
\(C_j(\lambda)\geq j+1\) for \(0\leq j\leq b\), by
using powers of the second variable alone. It follows that
\(D_{\mathrm{box}}\geq u+v+1=b\). In every variable split,
the part containing the second variable contributes at least
\(b+1\) monomials and the other part contributes at least one;
therefore \(T_b\geq b+1\). The hypothesis in
\eqref{eq:F015} consequently implies \(n\geq b\).

Let \(U\) have basis \(e_1,\ldots,e_r\), with dual variables
\(x_1,\ldots,x_r\), and put
\[
       \mu=(db-|\bar\lambda|,\lambda_2,\ldots,\lambda_r).
\]
By F-014 this is a partition and
\(a_\mu(d[b])=a_\lambda(d[n])\).
Let \(H_\mu\) be the complete highest-weight space of weight
\(\mu\) in \(\operatorname{Sym}^d\operatorname{Sym}^b U\).
Define the monomial subspace
\[
 W=\operatorname{span}\{e^\eta:|\eta|=b,\quad
                    \eta_i\leq\lambda_i\ (2\leq i\leq r)\}
                  \subseteq\operatorname{Sym}^b U.
\]
Every monomial in a polynomial of weight \(\mu\) has total
nonnegative weight \(\mu\). A generator \(e^\eta\) with
\(\eta_i>\mu_i=\lambda_i\) for some \(i\geq2\) cannot
appear in such a monomial. Therefore
\begin{equation}\label{eq:F015support}
                         H_\mu\subseteq\operatorname{Sym}^d W.
\end{equation}
This statement concerns every highest-weight vector, including
all multiplicity copies. It does not require \(W\) to be a
Borel-stable subspace.

Let \(\mathcal C\subseteq\operatorname{Sym}^b U^*\) be the
span of the dual monomials \(x^\eta\) with the same support
conditions. The differential pairing restricts to a nondegenerate
diagonal pairing between \(W\) and \(\mathcal C\), with
\(\langle e^\eta,x^\eta\rangle=\eta!\).
Consequently restriction of polynomial functions to
\(\mathcal C\) is injective on \(\operatorname{Sym}^d W\),
and in particular on \(H_\mu\).

Put \(p=n-b\). Multiplication
\(T:e^\eta\mapsto e_1^p e^\eta\) from
\(\operatorname{Sym}^b U\) to \(\operatorname{Sym}^n U\)
is upper-Borel equivariant up to the character \(g_{11}^p\).
Thus \(\operatorname{Sym}^d T\) takes \(H_\mu\) injectively
to highest-weight vectors of weight
\(\mu+dp e_1=\lambda\).
With \(U\) the span of the first \(r\) basis vectors of the
ambient space, these are also highest-weight vectors for the
ambient general linear group: raising operators involving a
later source index act by zero.
Consider actual form padding
\[
                 P:\mathcal C\longrightarrow
                       \operatorname{Sym}^n U^*,\qquad
                         f\longmapsto x_1^p f.
\]
As in O-012.b, the transpose \(R\) of padding is
\(\partial_{e_1}^p\), so
\[
                 RT(e^\eta)
                   =\frac{(\eta_1+p)!}{\eta_1!}e^\eta.
\]
This invertible diagonal operator preserves \(W\). Pulling
the lifted vectors back to \(P(\mathcal C)\) is consequently
the restriction of the invertible operator
\(\operatorname{Sym}^d(RT)\) on \(\operatorname{Sym}^d W\).
By \eqref{eq:F015support}, a basis of \(H_\mu\) remains
linearly independent after lifting and restriction.

It remains to put that restricted padded family in \(\Omega_n\).
The coefficient-box determinant construction O-014.a, obtained
by restricting both degree-split trees of O-012.a to the
downward-closed index sets above, gives
\[
           x_1^{n-b}\mathcal C\subseteq\Omega_n
                    \qquad\text{when }n\geq D_{\mathrm{box}}.
\]
Alternatively the bounded variable-split construction in M-009,
applied to degree \(b\), gives the same containment when
\(n\geq T_b\). In either case the displayed basis remains
linearly independent in the determinant coordinate ring.
Thus
\[
 K_n(\lambda)\geq\dim H_\mu
       =a_\mu(d[b])=a_\lambda(d[n]).
\]
The reverse inequality holds because that coordinate ring is
a quotient of the ambient polynomial ring. This proves
\eqref{eq:F015}.
\end{proof}

\paragraph{Exact scope of the improvement.}
The argument preserves all ambient highest-weight copies by
restricting to a linear space adapted to their support. It does
not replace the desired comparison by a positive occurrence
statement, and it does not infer containment of every padded
degree-\(b\) form. The rank criterion F-010 is now decided
whenever \eqref{eq:F015} holds: the determinant evaluation map
has its maximum possible rank. When the displayed inequality
fails, the construction supplies no rank inequality.
The permanent degree \(m\) enters only when applying separate
permanent-side shape constraints such as \(r\leq m^2\).

\paragraph{Grothendieck dictionary.}
The box is a finite torus-weight support in a locally free
symmetric-power module, with the grading/comodule interpretation
of SGA~3, I, 4.7.3. Its dual is an affine coefficient scheme.
The graph identities define integral morphisms from that scheme
to the determinant locus. The restriction and padding maps are
symmetric-power maps, and their differential transpose has
explicit factorial diagonal entries. Cauchy and the Schur
functor decomposition, Macdonald I(4.3) and Appendix~A(5.3)--(5.6),
identify the characteristic-zero highest-weight multiplicity
spaces. Invertibility of those factorials and semisimplicity
are asserted only in characteristic zero. This is a coherent
quotient calculation on the orbit-closure scheme; it does not
claim an effective Burnside lift or positivity of virtual
classes in unrelated fibers.
